\documentclass[10pt,a4paper]{amsart}
\usepackage[margin=19mm]{geometry}
\usepackage{amsmath,amssymb,mathtools}
\usepackage[scr=rsfs,cal=euler]{mathalfa}
\usepackage{xfrac}
\usepackage[unicode,hidelinks,bookmarks=false]{hyperref}
\newcommand{\ie}{i.e.\ }
\newcommand{\cf}{cf.\ }
\newcommand{\resp}{respectively\ }
\newcommand{\Subsection}{\subsection}
\providecommand{\nobreakdash}{\nobreak\hbox{-}}
\setlength{\emergencystretch}{2em}
\multlinegap0pt
\usepackage{amssymb}
\usepackage{mathtools}
\usepackage{xcolor}
\usepackage[shortlabels]{enumitem}
\usepackage{tikz}
\newtheorem{theorem}{Theorem}
\newtheorem{lemma}{Lemma}
\newtheorem{prop}{Proposition}
 \newcommand{\R}{\ensuremath{\mathbb{R}}}% real numbers
\newcommand{\holsp}[1][\eta]{\mathcal{C}^{#1}}
\newcommand{\kant}{\wass_1}
\newcommand{\pnorm}[1]{\left|#1\right|}
\newcommand{\shat}[1]{\vphantom{#1}\smash[t]{\widehat{#1}}}
\newcommand{\stilde}[1]{\vphantom{#1}\smash[t]{\widetilde{#1}}}
\newcommand{\wass}{\mathcal{W}}
\title{Rates of convergence in the multivariate weak invariance principle for nonuniformly hyperbolic maps}
\author{Nicholas Fleming-V\'azquez}
\date{Source: arXiv:2503.16358v2 (CC BY 4.0)}
\begin{document}
\maketitle

\noindent\emph{Source and licence.} Rates of convergence in the multivariate weak invariance principle for nonuniformly hyperbolic maps. Version arXiv:2503.16358v2 (CC BY 4.0), distributed by arXiv under the Creative Commons Attribution 4.0 International licence, http://creativecommons.org/licenses/by/4.0/, as recorded in the arXiv licence metadata for this exact version and shown on the abstract page https://arxiv.org/abs/2503.16358v2. Full licence notice: \texttt{licenses/arxiv-2503.16358-CC-BY-4.0.txt}.\par\smallskip

\noindent\textit{Editorial scope.} Counted unit: the article's theorem thm:rates\_wip\_lip (source0.tex lines 382--387). The article's complete original proof of it is delivered with it (source0.tex lines 683--722, one \texttt{proof} environment of 40 lines, which the article places away from the statement). No mathematics is written, repaired or re-derived: the statement and the proof are the article's own lines, in the article's own notation, and no retained line is substituted or re-flowed.  Retained uncounted as the source-local context the proof needs: lines 407--410; lines 495--498; lines 521--523; lines 553--555; lines 571--573; lines 629--632.  Nothing was omitted from any retained span, so the package carries no omission record.  No retained line contains a citation, so the wrapper carries no bibliography and no bibliography entry.  No retained line contains a figure environment, an image-inclusion command or drawing code, so no image is delivered and the package declares no asset dependency.  Editorial formatting only, in addition: an AMS \texttt{amsart} preamble that restates the article's own declarations and macros used by the retained lines, namely the environments \texttt{theorem}, \texttt{lemma}, \texttt{prop} on separate counters, together with the standard packages those lines need.  The retained environments print in this document as Theorem 1, Lemma 1, Proposition 1 in order of appearance, whereas the article prints the same items with the numbering of its own full document.  The complete native source of the article as distributed in the arXiv e-print for this exact version is bundled unchanged as \texttt{sources/175123/source0.tex}.
	\begin{lemma}\label{lem:wasserstein_interpolation}
	Let $1\le r< s$. Let $X_1,X_2$ be random elements of $D=D([0,1],\R^d)$ with laws in $ \mathcal{P}_r^s(D)$. Let $\theta = \frac{1/r-1/s}{1-1/s}$. Then 
	\[ \wass_r(X_1,X_2) \le \wass_1(X_1,X_2)^{\theta}(\wass_s(X_1,0)+\wass_s(X_2,0))^{1-\theta}. \]
\end{lemma}

	\begin{prop}\label{prop:maximal_ineq_birkhoff}
		There is a constant $ C>0 $ such that 
		$ \pnorm{\max_{1\le j\le n}|S_j|}_{2\gamma}\le Cn^{1/2}$ for all $n\ge 1$.
	\end{prop}

\begin{lemma}\label{lemma:blocking_step_error}There is a constant $ C>0 $ such that 
		$ \wass_{2\gamma}(W_n, \stilde W_n+R_n)\le Cn^{(1-a)\frac{1-\gamma}{2\gamma}} $ for all $ n\ge 1 $.
	\end{lemma}

\begin{lemma}\label{lem:small_blocks_negligible}
	There is a constant $C>0$ such that $\wass_{2\gamma}(\stilde W_n+R_n, \stilde W_n)\le Cn^{(b-a)/2}$ for all $n\ge 1$.
\end{lemma}

	\begin{lemma}\label{lemma:big_block_indep_approx}
		There is a constant $C>0$ such that $\kant(\stilde W_n, \shat W_n)\le Cn^{3/2-a-b\gamma}$ for all $n\ge 1$.
	\end{lemma}

	\begin{lemma}\label{lemma:approx_indep_big_blocks_bm}
		There is a constant $ C>0 $ such that for all $ n\ge 1 $,
		\[ \wass_{2\gamma}(\shat W_n,W) \le C(n^{a(1-\gamma)/2}+ n^{(b-a)/2}+n^{(1-a)\frac{1-\gamma}{2\gamma}}). \]
	\end{lemma}

\begin{theorem}\label{thm:rates_wip_lip}
	Let $ T:M\to M $ satisfy the Functional Correlation Bound with rate $ n^{-\gamma} $, $ \gamma>1 $, and suppose that $ v\in \holsp(M,\R^d)$ with mean zero. Let $1\le r<2\gamma$. Then there is a constant $ C>0 $ such that $ \wass_r(W_n,W)\le Cn^{-\kappa(\gamma,r)} $ for all $ n\ge 1 $, where 
	\begin{equation*}
		\kappa(\gamma,r)=\frac{(2\gamma-r)(2\gamma-1)(\gamma-1)}{2(8\gamma^3 - r(2\gamma^2 + 3\gamma - 1))}.
	\end{equation*}
\end{theorem}

\begin{proof}[Proof of Theorem~\ref{thm:rates_wip_lip}]
	By Lemmas~\ref{lemma:blocking_step_error} and~\ref{lem:small_blocks_negligible},
	\begin{align}
		\wass_{2\gamma}(W_n,\stilde W_n)&\le \wass_{2\gamma}(W_n, \stilde W_n+R_n)+\wass_{2\gamma}(\stilde W_n+R_n, \stilde W_n)\nonumber\\
		&\ll n^{(1-a)\frac{1-\gamma}{2\gamma}}+n^{(b-a)/2}.\label{eq:diff_with_big_blocks}
	\end{align}
	On the other hand, by Lemma~\ref{lemma:approx_indep_big_blocks_bm},
	\begin{align}
		\wass_{2\gamma}(\shat W_n,W)\ll n^{a(1-\gamma)/2}+ n^{(b-a)/2}+n^{(1-a)\frac{1-\gamma}{2\gamma}}.\label{eq:indep_blocks_vs_brownian}
	\end{align}
	Now by Proposition~\ref{prop:maximal_ineq_birkhoff}, $\wass_{2\gamma}(W_n,0)=O(1)$ and hence by~\eqref{eq:diff_with_big_blocks} we have that $\wass_{2\gamma}(\stilde W_n, 0)=O(1)$. Similarly,~\eqref{eq:indep_blocks_vs_brownian} implies that $\wass_{2\gamma}(\shat W_n, 0)=O(1)$. Hence by Lemma~\ref{lemma:big_block_indep_approx} and Lemma~\ref{lem:wasserstein_interpolation}, it follows that 
	\begin{align*}
		\wass_r(\stilde W_n, \shat W_n)\ll \wass_1(\stilde W_n, \shat W_n)^\frac{2\gamma-r}{r(2\gamma-1)}\ll n^{\frac{2\gamma-r}{r(2\gamma-1)}(3/2-a-b\gamma)}.
	\end{align*}
	By combining this with~\eqref{eq:diff_with_big_blocks} and~\eqref{eq:indep_blocks_vs_brownian}, we obtain that $ \wass_r(W_n, W)\ll n^{-\kappa(\gamma,r)} $ where
	\begin{align}
		\kappa(\gamma,r)=\min\Big\{(1-a)\frac{\gamma-1}{2\gamma},\, \frac{(2\gamma-r)(a+b\gamma-\frac{3}{2})}{r(2\gamma-1)},\, \frac{a-b}{2},\, \frac{a(\gamma-1)}{2}\Big\}.\label{eq:wip_implicit_Wp_bd}
	\end{align}
	Recall that $ 0<b<a<1 $ are arbitrary. Finally, we show how to choose $ a $ and $ b $ in order to obtain the expression for $ \kappa(\gamma,r) $ in the statement of this theorem. We choose $a$ and $b$ such that 
	\[ (1-a)\frac{\gamma-1}{2\gamma}=\frac{(2\gamma-r)(a+b\gamma-\frac{3}{2})}{r(2\gamma-1)}=\frac{a-b}{2}. \]
	Note that $ (1-a)\frac{\gamma-1}{2\gamma}=\frac{(2\gamma-r)(a+b\gamma-\frac{3}{2})}{r(2\gamma-1)}$ if and only if
	\begin{equation}\label{eq:b_eq_1}
		b = \frac{a(r(2\gamma+1)-4\gamma)+6\gamma-3r}{4\gamma^2-r}.
	\end{equation}
	On the other hand, $(1-a)\frac{\gamma-1}{2\gamma}=\frac{a-b}{2}$ if and only if 
	\begin{equation}\label{eq:b_eq_2}
		b = \frac{a(2\gamma-1)-\gamma+1}{\gamma}.
	\end{equation}
	By equating~\eqref{eq:b_eq_1} and~\eqref{eq:b_eq_2} and solving for $a$ we obtain that 
	\[ a=\frac{4\gamma^3+2\gamma^2 - 4r\gamma + r}{8\gamma^3 - r(2\gamma^2 + 3\gamma - 1)}\quad \text{ and }\quad (1-a)\frac{\gamma-1}{2\gamma}=\frac{(2\gamma-r)(2\gamma-1)(\gamma-1)}{2(8\gamma^3 - r(2\gamma^2 + 3\gamma - 1))}{.}
	\]
	Now since $r\le 2\gamma$, we have that $2\gamma^2-4r\gamma\ge -2\gamma^2-2r\gamma$ so
	\begin{align*}
		a \ge \frac{4\gamma^3-2\gamma^2 - 2r\gamma + r}{8\gamma^3 - r(2\gamma^2 + 3\gamma - 1)}=\frac{(2\gamma^2-r)(2\gamma-1)}{8\gamma^3 - r(2\gamma^2 + 3\gamma - 1)}
		\ge \frac{1-a}{\gamma}.
	\end{align*}
	It follows that
	\[ \kappa(\gamma,r)=\min\Big\{(1-a)\frac{\gamma-1}{2\gamma},\frac{a(\gamma-1)}{2} \Big\}=(1-a)\frac{\gamma-1}{2\gamma}=\frac{(2\gamma-r)(2\gamma-1)(\gamma-1)}{2(8\gamma^3 - r(2\gamma^2 + 3\gamma - 1))}.\]
	Since $\kappa(\gamma,r)>0$, by~\eqref{eq:wip_implicit_Wp_bd} we have that $0<b<a<1$, as required.
\end{proof}

\end{document}
